% GENERATED FILE. Edit canonical lessons, then run npm run generate:books.
% canonical-source-hash: fnv1a64:f2f080f919e59f0b
% canonical-lessons: HI-C546-sagun, HI-C546-asar, HI-C546-manhus, HI-C546-saubhagya, HI-C546-ittafaq

\chapter{Fate and Chance}
\label{ch:hi-fate-and-chance}
\hlchaptermodality{546}

\begin{chapteropening}
\textbf{By the end of this chapter:} \emph{I can talk about omens, bad luck and lucky chances.}

Complete the last lesson of chapter 546: I can talk about omens, bad luck and lucky chances.
\end{chapteropening}

\section[\texorpdfstring{\hi{शगुन} (śagun)}{शगुन (śagun)}]{\hi{शगुन} (śagun) — an omen (a good sign)}

\label{lesson:HI-C546-sagun}



\begin{quote}
Before the new one: say the Hindi for a piece (in chess or a board game), then the Hindi for defeat.
\end{quote}

\subsection*{You'll want to know: \hi{शगुन}}
\textbf{\hi{शगुन}} — \emph{śagun} — \textquotedblleft{}an omen (a good sign)\textquotedblright{}.

\begin{quote}
\hi{दादी} \hi{मानती} \hi{हैं} \hi{कि} \hi{यह} \hi{एक} \hi{अच्छा} \hi{शगुन} \hi{है।}
\end{quote}

\begin{quote}
\emph{Dādī māntī haiṁ ki yah ek acchā śagun hai.}
\end{quote}

\begin{quote}
Grandmother believes that this is a good omen.
\end{quote}

Masculine.

\begin{cousinweb}
\textbf{\hi{शगुन}} (\emph{śagun}) is probably a later borrowing of Sanskrit \emph{śakuna}, \textquotedblleft{}a bird\textquotedblright{}, which also meant \textquotedblleft{}an omen\textquotedblright{}, read from birds. The forms Hindi inherited, \emph{saun} and \emph{son}, \textquotedblleft{}a good omen\textquotedblright{}, are the ones CDIAL lists (CDIAL 12245); \emph{śagun} itself is not among them.
\end{cousinweb}

\begin{grammarlens}[title={What you've built}]
One more word for saying what happens if.
\end{grammarlens}

\subsection*{Your turn}
\begin{itemize}
\raggedright
  \item \emph{Say it:} \emph{śagun}
  \item \emph{Say it:} \emph{śagun}, once more
  \item \emph{Say it:} \emph{Dādī māntī haiṁ ki yah ek acchā śagun hai.} — Grandmother believes that this is a good omen.
  \item \emph{Recall:} say the Hindi for a piece (in chess or a board game), then the Hindi for defeat, then say \emph{śagun} again
\end{itemize}

\begin{tcolorbox}[breakable,colback=teal!4,colframe=teal!35!black,title={Before you move on}]
What does \hi{शगुन} mean? (\textquotedblleft{}An omen (a good sign)\textquotedblright{}.) Say it once more.
\end{tcolorbox}
\section[\texorpdfstring{\hi{आसार} (āsār)}{आसार (āsār)}]{\hi{आसार} (āsār) — signs; prospects (of something coming)}

\label{lesson:HI-C546-asar}



\begin{quote}
Before the new one: say the Hindi for defeat, then the Hindi for an omen (a good sign).
\end{quote}

\subsection*{You'll want to know: \hi{आसार}}
\textbf{\hi{आसार}} — \emph{āsār} — \textquotedblleft{}signs; prospects (of something coming)\textquotedblright{}.

\begin{quote}
\hi{आसमान} \hi{में} \hi{बादल} \hi{हैं}, \hi{बारिश} \hi{के} \hi{आसार} \hi{हैं।}
\end{quote}

\begin{quote}
\emph{Āsmān meṁ bādal haiṁ, bāriś ke āsār haiṁ.}
\end{quote}

\begin{quote}
There are clouds in the sky: rain looks likely.
\end{quote}

Masculine.

Used in the plural: \hi{बारिश} \hi{के} \hi{आसार} \hi{हैं}.

\begin{cousinweb}
\textbf{\hi{आसार}} (\emph{āsār}) is usually traced to Arabic \emph{āthār}, \textquotedblleft{}traces, signs\textquotedblright{}, through Persian.
\end{cousinweb}

\begin{grammarlens}[title={What you've built}]
One more word for saying what happens if.
\end{grammarlens}

\subsection*{Your turn}
\begin{itemize}
\raggedright
  \item \emph{Say it:} \emph{āsār}
  \item \emph{Say it:} \emph{āsār}, once more
  \item \emph{Say it:} \emph{Āsmān meṁ bādal haiṁ, bāriś ke āsār haiṁ.} — There are clouds in the sky: rain looks likely.
  \item \emph{Recall:} say the Hindi for defeat, then the Hindi for an omen (a good sign), then say \emph{āsār} again
\end{itemize}

\begin{tcolorbox}[breakable,colback=teal!4,colframe=teal!35!black,title={Before you move on}]
What does \hi{आसार} mean? (\textquotedblleft{}Signs; prospects (of something coming)\textquotedblright{}.) Say it once more.
\end{tcolorbox}
\section[\texorpdfstring{\hi{मनहूस} (manhūs)}{मनहूस (manhūs)}]{\hi{मनहूस} (manhūs) — unlucky; ill-omened}

\label{lesson:HI-C546-manhus}



\begin{quote}
Before the new one: say the Hindi for an omen (a good sign), then the Hindi for signs.
\end{quote}

\subsection*{You'll want to know: \hi{मनहूस}}
\textbf{\hi{मनहूस}} — \emph{manhūs} — \textquotedblleft{}unlucky; ill-omened\textquotedblright{}.

\begin{quote}
\hi{पुराने} \hi{लोग} \hi{कहते} \hi{थे} \hi{कि} \hi{टूटा} \hi{शीशा} \hi{मनहूस} \hi{होता} \hi{है।}
\end{quote}

\begin{quote}
\emph{Purāne log kahte the ki ṭūṭā śīśā manhūs hotā hai.}
\end{quote}

\begin{quote}
Old people used to say that a broken mirror is unlucky.
\end{quote}

\begin{cousinweb}
\textbf{\hi{मनहूस}} (\emph{manhūs}) is usually traced to Arabic \emph{manḥūs}, \textquotedblleft{}ill-omened\textquotedblright{}, through Persian.
\end{cousinweb}

\begin{grammarlens}[title={What you've built}]
One more word for saying what happens if.
\end{grammarlens}

\subsection*{Your turn}
\begin{itemize}
\raggedright
  \item \emph{Say it:} \emph{manhūs}
  \item \emph{Say it:} \emph{manhūs}, once more
  \item \emph{Say it:} \emph{Purāne log kahte the ki ṭūṭā śīśā manhūs hotā hai.} — Old people used to say that a broken mirror is unlucky.
  \item \emph{Recall:} say the Hindi for an omen (a good sign), then the Hindi for signs, then say \emph{manhūs} again
\end{itemize}

\begin{tcolorbox}[breakable,colback=teal!4,colframe=teal!35!black,title={Before you move on}]
What does \hi{मनहूस} mean? (\textquotedblleft{}Unlucky; ill-omened\textquotedblright{}.) Say it once more.
\end{tcolorbox}
\section[\texorpdfstring{\hi{सौभाग्य} (saubhāgya)}{सौभाग्य (saubhāgya)}]{\hi{सौभाग्य} (saubhāgya) — good fortune}

\label{lesson:HI-C546-saubhagya}



\begin{quote}
Before the new one: say the Hindi for signs, then the Hindi for unlucky.
\end{quote}

\subsection*{You'll want to know: \hi{सौभाग्य}}
\textbf{\hi{सौभाग्य}} — \emph{saubhāgya} — \textquotedblleft{}good fortune\textquotedblright{}.

\begin{quote}
\hi{आपसे} \hi{मिलना} \hi{मेरा} \hi{सौभाग्य} \hi{है।}
\end{quote}

\begin{quote}
\emph{Āpse milnā merā saubhāgya hai.}
\end{quote}

\begin{quote}
It is my good fortune to meet you.
\end{quote}

Masculine.

\begin{cousinweb}
\textbf{\hi{सौभाग्य}} (\emph{saubhāgya}) is a tatsama, taken whole from Sanskrit \emph{saubhāgya}, \textquotedblleft{}welfare, good fortune\textquotedblright{} (Monier-Williams).
\end{cousinweb}

\begin{grammarlens}[title={What you've built}]
One more word for saying what happens if.
\end{grammarlens}

\subsection*{Your turn}
\begin{itemize}
\raggedright
  \item \emph{Say it:} \emph{saubhāgya}
  \item \emph{Say it:} \emph{saubhāgya}, once more
  \item \emph{Say it:} \emph{Āpse milnā merā saubhāgya hai.} — It is my good fortune to meet you.
  \item \emph{Recall:} say the Hindi for signs, then the Hindi for unlucky, then say \emph{saubhāgya} again
\end{itemize}

\begin{tcolorbox}[breakable,colback=teal!4,colframe=teal!35!black,title={Before you move on}]
What does \hi{सौभाग्य} mean? (\textquotedblleft{}Good fortune\textquotedblright{}.) Say it once more.
\end{tcolorbox}
\section[\texorpdfstring{\hi{इत्तफ़ाक़} (ittafāq)}{इत्तफ़ाक़ (ittafāq)}]{\hi{इत्तफ़ाक़} (ittafāq) — chance; a chance happening}

\label{lesson:HI-C546-ittafaq}



\begin{quote}
Before the new one: say the Hindi for unlucky, then the Hindi for good fortune.
\end{quote}

\subsection*{You'll want to know: \hi{इत्तफ़ाक़}}
\textbf{\hi{इत्तफ़ाक़}} — \emph{ittafāq} — \textquotedblleft{}chance; a chance happening\textquotedblright{}.

\begin{quote}
\hi{इत्तफ़ाक़} \hi{से} \hi{हम} \hi{दोनों} \hi{एक} \hi{ही} \hi{ट्रेन} \hi{में} \hi{थे।}
\end{quote}

\begin{quote}
\emph{Ittafāq se ham donõ ek hī ṭren meṁ the.}
\end{quote}

\begin{quote}
By chance we were both on the same train.
\end{quote}

Masculine.

\begin{cousinweb}
\textbf{\hi{इत्तफ़ाक़}} (\emph{ittafāq}) is usually traced to Arabic \emph{ittifāq}, \textquotedblleft{}a coming together, agreement\textquotedblright{}, through Persian.
\end{cousinweb}

\begin{grammarlens}[title={What you've built}]
One more word for saying what happens if.
\end{grammarlens}

\subsection*{Your turn}
\begin{itemize}
\raggedright
  \item \emph{Say it:} \emph{ittafāq}
  \item \emph{Say it:} \emph{ittafāq}, once more
  \item \emph{Say it:} \emph{Ittafāq se ham donõ ek hī ṭren meṁ the.} — By chance we were both on the same train.
  \item \emph{Recall:} say the Hindi for unlucky, then the Hindi for good fortune, then say \emph{ittafāq} again
\end{itemize}

\begin{tcolorbox}[breakable,colback=teal!4,colframe=teal!35!black,title={Before you move on}]
What does \hi{इत्तफ़ाक़} mean? (\textquotedblleft{}Chance; a chance happening\textquotedblright{}.) Say it once more.
\end{tcolorbox}
